% =====================================================================
%  2 -- CONSERVATION AS LINEAR ALGEBRA  (source node S7).
%
%  Source: tier2.md 0.4 (743-837, RE-HOMED here from Part 0 -- it is the
%          bridge from the astrophysics to the algebra, and every one of
%          its three subsections is a statement about nu), then Part I
%          in full and in order: header (1048-1058), I.1 (1062-1126),
%          I.2 (1130-1272), I.2b (1276-1398), I.3 (1402-1437),
%          I.4 (1441-1491), I.5 (1495-1533), I.6 (1537-1572),
%          I.7 (1576-1611), I.7b (1615-1694), I.8 (1698-1763),
%          I.9 (1767-1809), I.10 (1813-1858), I.10b (1862-2007),
%          I.11 (2011-2051), I.12 (2055-2103), I.13 (2107-2163),
%          I.14 self-check (2167-2185).
%  Nothing from this range is re-homed out.
% =====================================================================

\section{Conservation as linear algebra}
\label{part:conservation}

\begin{repopointer}
\textbf{Files:} \code{graph/stoich.py}, \code{graph/projector.py},
\code{graph/isotopes.py}, \code{graph/network.py}, \code{graph/export.py},
\code{graph/metrics.py}. \quad
\textbf{Oracles:} \code{tests/test\_conservation.py},
\code{tests/test\_projector.py} --- \S\ref{sec:oracles7}. \quad
\textbf{Notebook:} \code{03-graph-conservation.ipynb}.
\end{repopointer}

\textbf{The floor.} \code{tests/test\_conservation.py} is the one test in the
repo that is allowed to block training. This section is what it means.

The plan of attack: state the theorem that makes conservation structural rather
than learned (\S\ref{sec:whystructural}), build $\nu$ (\S\ref{sec:numap}),
prove that conservation laws are its left null space and \emph{measure} that
null space (\S\ref{sec:leftnullspace}), place the whole object in the two
literatures that already own it (\S\ref{sec:crn}), watch the weak sector break
charge conservation (\S\ref{sec:weakbreaks}), repair it with an independent
lepton ledger (\S\ref{sec:ledgers}--\S\ref{sec:betaplus}), assemble $C$
(\S\ref{sec:cmatrix}), get Target A's theorem for free
(\S\ref{sec:targetatheorem}) and note what it does not give
(\S\ref{sec:positivity}), derive the Target B projector twice
(\S\ref{sec:projector}--\S\ref{sec:laundering}), prove the two targets reach the
same set (\S\ref{sec:equivalence}), measure how big that set actually is
(\S\ref{sec:effdim}), and read the graph, the code and the oracle
(\S\ref{sec:graphK}--\S\ref{sec:oracles7}).

% =====================================================================
\subsection{Why conservation is structural, not a penalty term}
\label{sec:whystructural}

\subsubsection{Conservation laws are left null vectors}
\label{sec:leftnull}

Here is the theorem the entire S7 node rests on, stated cleanly.

Let $\nu \in \mathbb{R}^{n\times r}$ be the stoichiometric matrix. A vector
$c \in \mathbb{R}^n$ defines a \textbf{conserved quantity} $Q = c\cdot Y$ of the
dynamics $\dot y = \nu R$ \textbf{iff} $c$ is a left null vector of $\nu$:
%
\begin{equation}
\frac{\dd}{\dd t}(c\cdot Y) \;=\; c\Tr\nu R \;=\; 0 \ \ \forall R
\quad\Longleftrightarrow\quad c\Tr\nu = 0.
\label{eq:leftnull}
\end{equation}
%
The ``for all $R$'' is the crucial quantifier. It says the conservation law
holds \emph{regardless of the rates} --- regardless of temperature, density,
screening, whether the rate library is right, whether the network is converged,
or whether the thing producing $R$ is a stiff ODE solver or an untrained neural
network with random weights. Conservation is a property of the
\textbf{stoichiometry}, and stoichiometry is integer bookkeeping.

The set of conserved quantities is therefore exactly the left null space of
$\nu$, a linear subspace, and its dimension is $n - \rank(\nu)$. This is
measurable, and \S\ref{sec:rank} measures it.

% ---------------------------------------------------------------------
\subsubsection{The soft-penalty failure mode, quantified}
\label{sec:softpenalty}

Suppose instead you train with a penalty: minimise
$L = L_{\rm data} + \lambda\|C\cdot \dd Y\|^2$. Three things go wrong, in
increasing order of severity.

\paragraph{(i) The violation is $O(1/\lambda)$, never zero.} At the optimum the
penalty gradient balances the data gradient:
$\|C\cdot \dd Y\| \sim \|\nabla L_{\rm data}\| / (2\lambda\|C\|^2)$. You buy
accuracy in conservation with accuracy in the fit, and never reach exact.

\paragraph{(ii) The violation is systematic.} This is the part people miss. The
penalty is minimised \emph{on the training distribution}; the residual violation
is whatever the data loss pays for. That is a \emph{learned, input-dependent}
function, not noise --- it has the same sign for similar inputs. Over a rollout
on a smooth trajectory, consecutive states are similar, so the violations
\emph{add coherently}. Referring to \S\ref{sec:budget}'s table: a soft-penalty
violation accumulates as $N\cdot\delta$, not $\sqrt N\cdot\delta$.

Do the arithmetic against the right budget, because it is easy to compare an
\emph{accumulated} violation with a \emph{per-step} tolerance and get a wrong
answer by three orders of magnitude. With $N = 1.6\times10^{3}$ the two budgets
are $\delta \le 3\times10^{-6}$ per step and $5\times10^{-3}$ over the
trajectory (\S\ref{sec:budget}):

\begin{center}
\begin{tabularx}{\textwidth}{@{}L{3.0cm} L{4.2cm} Y@{}}
\toprule
\textbf{per-step violation} & \textbf{accumulated over
  $N = 1.6\times10^{3}$} & \textbf{as a share of the $5\times10^{-3}$ total
  budget} \\
\midrule
$10^{-8}$ & $1.6\times10^{-5}$ & 0.3\% --- genuinely negligible \\
$10^{-7}$ & $1.6\times10^{-4}$ & 3\% \\
$\mathbf{10^{-6}}$ & $\mathbf{1.6\times10^{-3}}$ &
  \textbf{32\% --- a third of the budget, on bookkeeping} \\
$3\times10^{-6}$ & $4.8\times10^{-3}$ &
  96\% --- the entire budget, before the model errs at all \\
\bottomrule
\end{tabularx}
\end{center}

So the honest statement is \emph{not} that a penalty reaching $10^{-8}$ is
disqualifying: at $10^{-8}$ it is fine. It is that (a) the penalty gives you no
control over where on that table you land --- \S\ref{sec:softpenalty}(i) says
the residual is set by $\lambda$ against the data gradient, not by a
specification --- and (b) because the violation is systematic, the column that
matters is the \emph{accumulated} one, so every decade you fail to buy costs a
factor of ten on a budget you are also spending on real modelling error. The
exact route costs nothing and lands at row zero.

\paragraph{(iii) It corrupts the physics you care about.} The cheapest way for
an optimiser to reduce $\|C\cdot \dd Y\|$ is to shrink $\dd Y$. The charge row
of $C$ includes the weak-sector $\dYe$ signal. A penalty therefore applies
gradient pressure \emph{against} the one quantity the project exists to predict.
This is \code{CLAUDE.md} invariant \#3 stated as an optimisation fact: ``a
design that zeroes it to make the drift residuals vanish is wrong.''

\begin{invariant}
\textbf{The alternative costs nothing.} $\dd Y = \nu\netflux$ is a matrix
multiply with a fixed, integer-valued, precomputed $\nu$. It is \emph{cheaper}
than the penalty (no extra loss term, no $\lambda$ to tune) and exactly
conserving for any \netflux. There is no trade-off being made here --- which is
why the repo treats it as a floor rather than a feature.
\end{invariant}

% ---------------------------------------------------------------------
\subsubsection{The reachable manifold, restated}
\label{sec:reachable}

Tier~0 \S IV.2 introduced the reachable manifold: the set of compositions that a
one-zone burner can actually be in. Tier~2 gives the \emph{linear-algebraic}
version of the same idea, and it is sharper.

At any state $Y$, the set of \emph{permissible} increments is
%
\begin{equation}
\mathcal{M} \;=\; \bigl\{\,\dd\Yt \in \mathbb{R}^{n+3} :
                   C\,\dd\Yt = 0\,\bigr\},
\label{eq:manifold}
\end{equation}
%
a linear subspace of codimension $\rank(C) = 3$. Target A's outputs live in
$\image(\nut) \subseteq \mathcal{M}$ automatically. Target B's outputs are
forced into $\mathcal{M}$ by projection. \S\ref{sec:equivalence} proves those
two sets are \textbf{the same} --- which reframes the Target A/B choice
entirely: it is not about \emph{what can be expressed}, it is about
parameterisation and conditioning.

% =====================================================================
\subsection{\texorpdfstring{$\nu$}{nu} as a linear map}
\label{sec:numap}

\code{graph/stoich.py} builds $\nu$ in nine lines:

\begin{srclisting}
for j, rate in enumerate(rates):
    for sp in rate.reactants:
        nu[idx[sp], j] -= 1.0
    for sp in rate.products:
        nu[idx[sp], j] += 1.0
\end{srclisting}

That is the entire construction. $\nu[i, j]$ is the \textbf{net} change in the
count of species $i$ per unit progress of reaction $j$: products positive,
reactants negative, with repeated participants accumulating (triple-$\alpha$
gives $\nu[\code{he4}, j] = -3$).

Three structural facts follow immediately, and each matters later.

\paragraph{(a) Column order is the contract.} $\nu$'s columns are
\code{list(rc.get\_rates())} in pynucastro's order, and \emph{every} downstream
array --- $\lambda$, $R$, \fp, \netflux, \kap, \Phifl, \code{pair\_col},
\code{weak\_mask}, $Q$, \code{chapter} --- is indexed by the same $j$.
\code{fluxes/compile.py} asserts this alignment against \code{Stoich.rate\_fnames}
and raises if a v-flag replacement changed a column identity.
\foreignreg{1}{R-17}'s canonical key exists so that ``same reaction'' is
decidable; this is where that decidability is cashed in. A column-order bug
would be silent and catastrophic: every conservation test would still pass
(\S\ref{sec:leftnullspace} holds for \emph{any} $\nu$), while every flux would
be attached to the wrong reaction.

\paragraph{(b) $\nu$ is \emph{net}, so catalysts vanish.} A species that appears
on both sides (consumed and produced) gets the difference. This is why
\code{export.bipartite\_digraph} adds \emph{both} an I$\to$R and an R$\to$I edge
for such a species: the graph keeps the information $\nu$ throws away. For the
conservation algebra the net form is what you want; for message passing it is
not, and the graph is the compensating structure.

\paragraph{(c) $\nu$ is extremely sparse and hub-dominated.}\derivedhere{}
Measured on the reconciled export:

\begin{center}
\begin{tabularx}{\textwidth}{@{}Y L{3.4cm} L{3.4cm}@{}}
\toprule
 & \textbf{\msn{80}} & \textbf{\msn{151}} \\
\midrule
shape ($n \times r$) & $80 \times 607$ & $151 \times 1518$ \\
nonzeros & 2012 & 5004 \\
density & 4.1\% & 2.2\% \\
participants per column (mean\,/\,max) & 3.31\,/\,5 & 3.30\,/\,5 \\
\nuc{He}{4} appears in & 319 columns (52.6\%) & 695 (45.8\%) \\
\nuc{H}{1}\,/\,$n$ appear in & 246\,/\,247 & 644\,/\,643 \\
next-highest species degree & \nuc{C}{12}, \nuc{O}{16}: 32 &
  \nuc{O}{16}: 35 \\
\bottomrule
\end{tabularx}
\end{center}

\begin{srclisting}
import numpy as np
d = np.load("data/stoich/nu_mesa80.npz"); nu = d["nu"]; sp = list(d["species"])
part = nu != 0
print(part.sum(), part.sum()/nu.size, part.sum(0).mean(), part.sum(0).max())
print([(sp[i], int(part.sum(1)[i])) for i in np.argsort(-part.sum(1))[:6]])
\end{srclisting}

The degree distribution has a cliff: three light species with degree 250--700,
then everything else at $\le 35$. \textbf{\nuc{He}{4}, \nuc{H}{1} and $n$ are
not species in this network so much as they are the network's wiring.} This
single fact explains \S\ref{sec:twoclusters}'s reservoir picture, the graph
radius of 1 in the isotope projection (\S\ref{sec:graphK}), and a real
architectural hazard: a message-passing model can route almost any influence
through three hub nodes in two hops, which is both why $K = 5$ suffices and why
an ablation that damages the hub features would be uninformative about whether
the model learned the chemistry.

% =====================================================================
\subsection{Conservation is the left null space --- with the measurement}
\label{sec:leftnullspace}

\S\ref{sec:leftnull} stated the theorem. Here it is used.

\subsubsection{The rank measurement --- and why it is the most informative
               number in S7}
\label{sec:rank}

The left null space of $\nu$ is $\operatorname{null}(\nu\Tr)$ --- the set of all
conserved linear combinations. Its dimension is $n - \rank(\nu)$. Measure it:

\begin{srclisting}
import numpy as np
d  = np.load("data/stoich/nu_mesa80.npz")
nu, A, Z = d["nu"], d["A"], d["Z"]
U, S, Vt = np.linalg.svd(nu)
k = (S > 1e-12 * S[0]).sum()                 # numerical rank
L = U[:, k:]                                 # basis of the left null space
print(nu.shape, k, nu.shape[0] - k, np.linalg.norm(L.T @ (A/np.linalg.norm(A))))
\end{srclisting}

\derivedhere, at the reconciled export:

\begin{center}
\begin{tabularx}{\textwidth}{@{}Y L{2.6cm} L{2.6cm}@{}}
\toprule
\textbf{quantity} & \textbf{\msn{80}} & \textbf{\msn{151}} \\
\midrule
$\rank(\nu)$ & 79 & 150 \\
\textbf{left-nullity of $\nu$} & \textbf{1} & \textbf{1} \\
$\|$proj of $\hat A$ onto left-null$(\nu)\|$ & 1.0000 & 1.0000 \\
left-nullity of $\nu$ restricted to \emph{strong\,/\,EM} columns &
  \textbf{2} & \textbf{2} \\
\ldots and it is exactly $\mathrm{span}\{A, Z\}$ (both principal angles 0) &
  \yes & \yes \\
$\rank(\nut)$ (extended, with lepton rows) & 80 & 151 \\
\textbf{left-nullity of $\nut$} & \textbf{3} & \textbf{3} \\
\ldots and it is exactly $\mathrm{rowspace}(C)$ (all three principal angles 0) &
  \yes & \yes \\
right-nullity of $\nu$ ($= r - \rank$) & \textbf{528} & \textbf{1368} \\
$\cond(\nu)$ (extreme nonzero singular values) & 41.57 & 57.86 \\
$\cond(\nut)$ & 40.90 & 54.54 \\
$\cond(CC\Tr)$ & $3.947\times10^{4}$ & $1.029\times10^{5}$ \\
\bottomrule
\end{tabularx}
\end{center}

Read those rows one at a time; each is a physical statement.

\paragraph{Left-nullity of $\nu$ is 1, and the null vector is $A$.} Mass number
is the \emph{only} linear combination of nuclear abundances conserved by every
column. There are no accidental extra conservation laws hiding in the network
--- no coincidental degeneracy that a naive projector would have to be told
about, and no missing constraint the gate could be blind to. It also means the
constraint set cannot be \emph{over}-specified: $C$'s three rows are all doing
work.

\paragraph{Restricted to strong\,/\,EM columns, the nullity is 2, and the space
is exactly $\mathrm{span}\{A, Z\}$.} This is \S\ref{sec:weakbreaks}'s fact,
measured: charge is conserved by the strong sector and only by the strong
sector. The measurement (principal angles between $\mathrm{span}\{A, Z\}$ and
the computed left null space, both exactly 1.0 in cosine) is a much stronger
statement than ``$Z\cdot\nu = 0$ on strong columns'' --- it says there is
nothing \emph{else} conserved there either.

\begin{result}
\textbf{Corollary, and it is the most consequential thing in this table.} Every
increment the strong sector can produce satisfies $Z\cdot \dd Y = 0$ exactly,
i.e. $\dYe = 0$. The strong sector is fast and carries essentially all of the
composition change; therefore essentially all of the \emph{composition variance}
is \Ye-neutral, and the \Ye{} signal is structurally confined to the slow, small
weak sector. Measured consequence: over the visited compositions, $\Var(X)$
converges to 99.7\%\,/\,98.1\% by 30 principal components while $\Var(\Ye)$
stalls at 93.0\%\,/\,95.3\% --- the \Ye{} signal is smeared across a long
low-variance tail, and the per-step budget needs $1 - 6.5\times10^{-8}$ of it.
Variance-truncated compression of the composition is therefore arithmetically
excluded, not merely risky. \S\ref{sec:effdim} measures it; this row is the
reason it must come out that way.
\end{result}

\paragraph{Left-nullity of $\nut$ is 3, and the space is exactly
$\mathrm{rowspace}(C)$.} This is the completeness statement, and it is the one I
would not have guessed. It says $C$ is neither missing a constraint (else the
nullity would exceed 3) nor imposing a redundant one (else $\rank(C) < 3$,
caught by \code{build\_projector}). \textbf{The projector removes exactly the
physical constraints and nothing else.}

\begin{srclisting}
nx, C = d["nu_ext"], d["C"]
Ue, Se, _ = np.linalg.svd(nx); ke = (Se > 1e-12*Se[0]).sum()
Qc, _ = np.linalg.qr(C.T)
print(np.linalg.svd(Qc.T @ Ue[:, ke:], compute_uv=False))   # -> [1. 1. 1.]
\end{srclisting}

\paragraph{Right-nullity is 528 (1368).} The map $\netflux \mapsto \nu\netflux$
has a 528-dimensional kernel. Take this seriously: \textbf{the flux vector is
not determined by the abundance change.} Given a target $\Delta Y$ there is a
528-parameter family of \netflux{} producing it. Consequences, all of which
reappear later:

\begin{itemize}
\item A loss defined only on $\Delta Y$ leaves 528 directions of \netflux{}
  completely unconstrained. Target A therefore \emph{needs} flux-space
  supervision, which is why \code{integrate.py} exists
  (\S\ref{sec:augmented}--\S\ref{sec:phinotrecoverable}) and why the shipped
  Zenodo labels (which carry only $\Delta X$) are insufficient for it.
\item Conversely, the redundancy is what makes the mask meaningful: killing a
  column changes \netflux{} without necessarily changing $\nu\netflux$ much.
\item And it is why $\cond(\nu)$ alone does not settle the kill-test.
  $\cond(\nu) = 41.6$ is benign; the kill-test's $\cond(\Sact)$ restricts to the
  \emph{active} columns, which is a different matrix (\resultsref{2026-07-12}:
  41.8\,/\,57.4, essentially the full-$\nu$ value, because the active set turned
  out to be everything).
\end{itemize}

\paragraph{$\cond(\nu) \approx 42$--58 is small, and that is a real result.} It
bounds how much a relative error in \netflux{} can be amplified into $\dd Y$:
$\|\delta(\nu\netflux)\|/\|\nu\netflux\| \le
\cond(\nu)\cdot\|\delta\netflux\|/\|\netflux\|$ for \netflux{} in the row space.
Two orders of magnitude of headroom below the $10^{6}$ gate. The dangerous
conditioning in this project is \emph{not} in $\nu$; it is in the cancellation
within \netflux{} itself (\S\ref{sec:kappa}).

% ---------------------------------------------------------------------
\subsubsection{Why the tests can demand exactly 0.0}
\label{sec:exactzero}

\code{tests/test\_conservation.py::test\_exported\_columns\_conserve\_exactly}
asserts

\begin{srclisting}
assert dA[j] == 0.0
\end{srclisting}

Zero tolerance. That is unusual enough to justify.

$\nu$'s entries are small integers --- measured over both exports, the set is
exactly $\{-3, -2, -1,$ $+1, +2, +3\}$, with max 5 participants per column and no
coefficient beyond 3 ($\pm3$ is triple-$\alpha$ and its reverse). $A$ and $Z$
are integers $\le 151$ and $\le 28$. The column sum $\sum_i A_i\nu_{ij}$ is
therefore a sum of at most 5 products of integers bounded by ${\sim}450$ in
magnitude. \textbf{Every intermediate value is exactly representable in float64}
(integers up to $2^{53}$), and IEEE-754 addition of exactly-representable
integers whose exact sum is also representable is exact \citep{ieee754,
goldberg1991}. So the computed sum is the exact mathematical sum, and the exact
sum is 0 by construction of the reaction.

This is why \code{build\_stoich} uses \code{!= 0.0} rather than a tolerance in
its fail-loud validation, and why \code{graph/metrics.drift\_metrics} reports
$D_A = D_Q = 0.0$ exactly rather than ``$\le 10^{-16}$''. \resultsref{2026-07-08}
records the column drifts as exactly 0.0.

The \emph{random-\netflux} drift test is a different animal, because \netflux{}
is not an integer. There, $\sum_i A_i(\nu\netflux)_i =
\sum_j (\sum_i A_i\nu_{ij})\netflux_j = \sum_j 0\cdot\netflux_j$ mathematically,
but the computation does the sums in the other order and float64 rounding of the
intermediate $\nu\netflux$ survives. The residual is bounded by float64
$\varepsilon$ times the \emph{gross} throughput
$G = |A|\cdot(|\nu|\cdot|\netflux|)$, which is exactly the tolerance the module
docstring defines:
%
\begin{equation}
\text{tol} \;=\; \max\!\bigl(10^{-12} s,\ 10^{-13} G\bigr),
\qquad s = \min\bigl(1, \max_j|\netflux_j|\bigr).
\label{eq:drifttol}
\end{equation}
%
The rationale is Tier~0 \S V.4's: an absolute bound is the right statement at
$O(1)$ magnitudes and \emph{vacuous} at $\netflux \sim 10^{-20}$ (where any
implementation passes), so the small end gets a relative-to-gross bound instead
--- at $10^{-13}\cdot G$, about a thousand times looser than float64
$\varepsilon$ accumulation, i.e. tight enough to catch a real bug and loose
enough not to be a fp-noise flake. \textbf{Every threshold in this project sits
just above a measured floor} (Tier~1's closing habit); this one sits above
$\varepsilon\cdot G$.

% =====================================================================
\subsection{The same object in another literature --- CRN theory, deficiency,
            and the flux cone}
\label{sec:crn}

Everything \S\ref{sec:leftnullspace} just derived --- $\nu$, its left null space
as conservation laws, its right null space as flux redundancy, the reachable set
as an affine subspace intersected with the positive orthant --- is standard
material in two mature literatures that nuclear astrophysics and this repo do
not appear to cite. Worth knowing here rather than later: the vocabulary unlocks
fifty years of results, and one of the theorems very nearly applies.

\begin{itemize}
\item \textbf{Chemical Reaction Network Theory} \citep{hornjackson1972};
  Feinberg's deficiency theory, 1972--1987 \citep{feinberg1972, feinberg1987},
  with \citet{horn1972} as co-founder --- collected in \citet{feinberg2019};
\item \textbf{Stoichiometric\,/\,metabolic network analysis}
  \citep{schuster1994} on elementary flux modes; flux balance analysis in the
  Palsson school \citep{varma1994, orth2010}; extreme pathways are
  \citet{schilling2000}, with \citet{klamt2003} the EFM/EP comparison.
\end{itemize}

\subsubsection{The vocabulary map}
\label{sec:crnvocab}

Worth memorising, because it unlocks fifty years of results and algorithms:

\begin{center}
\begin{tabularx}{\textwidth}{@{}Y Y L{3.6cm}@{}}
\toprule
\textbf{this project} & \textbf{CRN\,/\,systems biology} &
\textbf{where in Tier~2} \\
\midrule
$\nu$, stoichiometric matrix & $S$, the stoichiometric matrix &
  \S\ref{sec:numap} \\
left null space of $\nu$ & \emph{conservation relations}, moiety conservation &
  \S\ref{sec:rank} \\
right null space of $\nu$ & the \textbf{null space of $S$}; steady-state flux
  space & \S\ref{sec:rank}, \S\ref{sec:phinotrecoverable} \\
\netflux, net flux vector & $v$, the flux vector & \S\ref{sec:lambdatoR} \\
$Y_0 + \image(\nu) \cap \{Y \ge 0\}$ &
  \textbf{stoichiometric compatibility class} &
  Tier~0 \S IV.2, \S\ref{sec:positivity} \\
the equilibrium mask & \emph{blocked}\,/\,\emph{fixed} reactions &
  \S\ref{sec:departure} \\
bridge reactions carrying the net flow & \textbf{elementary flux modes},
  extreme pathways & \S\ref{sec:twoclusters} \\
Target A (predict \netflux, map through $\nu$) & flux-based network
  parameterisation & \S\ref{sec:targetatheorem} \\
\bottomrule
\end{tabularx}
\end{center}

\subsubsection{Deficiency --- measured, and the honest answer}
\label{sec:deficiency}

The headline theorems of CRN theory are stated in terms of the
\textbf{deficiency}
%
\begin{equation}
\delta \;=\; n_{\mathrm{complexes}} \;-\; \ell \;-\; \rank(\nu),
\label{eq:deficiency}
\end{equation}
%
with $n_{\rm complexes}$ the number of distinct reactant/product multisets and
$\ell$ the number of linkage classes (connected components of the complex
graph). The \textbf{Deficiency Zero Theorem} \citep{feinberg1987} says: a weakly
reversible network with $\delta = 0$ has, under mass-action kinetics, for
\emph{every} choice of positive rate constants, exactly one positive equilibrium
in each positive stoichiometric compatibility class, and it is asymptotically
stable relative to its compatibility class. That would be an extraordinarily
strong statement about this system --- existence and uniqueness of NSE, and
stability of the QSE manifold, \emph{for free}, independent of the rate library.

It does not apply. \derivedhere:

\begin{srclisting}
# complexes = distinct reactant / product multisets, one edge per column
\end{srclisting}

\begin{center}
\begin{tabularx}{\textwidth}{@{}Y L{2.6cm} L{2.6cm}@{}}
\toprule
 & \textbf{\msn{80}} & \textbf{\msn{151}} \\
\midrule
complexes $n_c$ & 278 & 561 \\
linkage classes $\ell$ & 73 & 84 \\
$\rank(\nu)$ & 79 & 150 \\
\textbf{deficiency $\delta$} & \textbf{126} & \textbf{327} \\
strongly connected components (directed complex graph) & 81 & 93 \\
weakly connected components & 73 & 84 \\
\textbf{weakly reversible?} & \textbf{no} & \textbf{no} \\
\bottomrule
\end{tabularx}
\end{center}

Both hypotheses fail, and instructively:

\begin{itemize}
\item \textbf{$\delta = 126$\,/\,327 is enormous}, not marginal. Deficiency
  counts the mismatch between the \emph{complex}-level and \emph{species}-level
  descriptions; a network of 607 reactions built from three light particles
  interacting with 77 heavy species has an enormous number of distinct complexes
  relative to its stoichiometric rank. No amount of tidying reduces this to
  zero.
\item \textbf{Weak reversibility fails by exactly the weak sector.} The
  SCC\,/\,WCC gap is 8 and 9 --- the same order as the number of one-way weak
  channel families. $\beta$-decay and electron capture have no inverse in this
  regime (\S\ref{sec:pairmap}), which is a \emph{physical} irreversibility, not
  a modelling omission.
\end{itemize}

So: the theorems give nothing. That is a real, checkable, negative result, and
it is worth having because the natural hope (``surely CRN theory tells us the
QSE manifold is stable'') is now closed rather than open.

\subsubsection{What CRN theory \emph{does} give: the flux cone}
\label{sec:fluxcone}

The part that transfers directly is the geometry of the flux space.

With irreversibility constraints (weak columns are one-way, $\netflux_j \ge 0$),
the set of admissible steady-state flux vectors is a \textbf{polyhedral cone}
%
\begin{equation}
\mathcal{C} \;=\; \{\netflux : \nu\netflux = 0,\
                   \netflux_j \ge 0 \ \forall j \in \text{irrev}\},
\label{eq:fluxcone}
\end{equation}
%
and the \textbf{elementary flux modes} are its support-minimal flows --- the
minimal sets of reactions that can sustain a steady flow; they coincide with the
cone's extreme rays once reversible reactions are split into forward/backward
halves \citep{gagneur2004}. ``Which small set of reactions carries the net flow
from the silicon group to the iron group'' (\S\ref{sec:twoclusters}, the whole
bridge programme) \emph{is} the extreme-pathway question, and it has:

\begin{itemize}
\item an exact algorithm (the double-description method, \citealp{motzkin1953};
  \citealp{fukuda1996}; \code{efmtool} --- \citealp{terzer2008} --- and
  successors),
\item a large literature on \textbf{network reduction} driven by the same
  question,
\item and known scaling limits (EFM enumeration is combinatorially explosive ---
  \citealp{klamt2002} --- which is why flux-\emph{sampling} and flux-balance
  optimisation are used instead at scale).
\end{itemize}

The repo's bridge discovery (\code{scripts/step5\_bridges.py}) is an empirical,
flux-weighted, per-state version: rank columns by $|\nu\netflux|$ across the
group boundary. That is a perfectly reasonable estimator, and it answers a
slightly different question --- \emph{which bridges are used on this manifold},
not \emph{which bridges exist}. Both are worth having, and the second is
currently unasked.

\textbf{What remains open} (register: \reg{R-09}). Not because the project is
doing anything wrong, but because (i) a reviewer from computational systems
biology will ask why elementary flux modes are not cited, (ii) the reduction
literature there is directly relevant to the network-reduction prior art
\code{tier1.md} \S VIII.C.7 already flags as under-surveyed, and (iii) the flux
cone is the correct geometric object for ``what fluxes are admissible'' ---
exactly what a flux-predicting emulator should be constrained to. A literature
and framing gap, not a correctness gap.

% =====================================================================
\subsection{The weak sector breaks charge conservation --- on purpose}
\label{sec:weakbreaks}

For a strong or electromagnetic reaction, both $A$ and $Z$ are conserved by the
\emph{nuclei alone}:
%
\begin{equation}
\sum_i A_i \nu_{ij} = 0, \qquad
\sum_i Z_i \nu_{ij} = 0 \qquad (j \text{ strong\,/\,EM}).
\label{eq:strongcons}
\end{equation}
%
For a weak reaction, the first still holds --- $\beta$-decay and electron
capture move a nucleon between the proton and neutron states without changing
$A$ --- but the second does not:
%
\begin{equation}
\sum_i Z_i \nu_{ij} = \Delta Z_j = \pm 1 \qquad (j \text{ weak}).
\label{eq:weakdz}
\end{equation}
%
Measured over the whole export: \code{(Z @ nu)[\textasciitilde weak]} is exactly
0 everywhere, and \code{(Z @ nu)[weak]} takes exactly the values $\{-1, +1\}$
--- no other value occurs, in either network.\derivedhere

This is not a defect to be repaired; \textbf{it is the physics the project
exists to predict.} $\Ye = \sum_i Z_iY_i$, so
%
\begin{equation}
\frac{\dd \Ye}{\dd \netflux_j} \;=\; \sum_i Z_i \nu_{ij} \;=\; \Delta Z_j,
\label{eq:dyedphi}
\end{equation}
%
and a network in which $\sum_i Z_i\nu_{ij} = 0$ for \emph{every} column is a
network in which \Ye{} can never change. \code{CLAUDE.md} invariant \#3 is the
guard against ``fixing'' it: any design that zeroes the weak \dYe{} to make
residuals vanish has deleted the answer.

So the constraint cannot be ``charge is conserved''. It has to be \textbf{charge
is conserved once you count the leptons}, which requires knowing, per column,
where the charge went. That is the ledger.

% =====================================================================
\subsection{Lepton ledgers --- and the tautology trap}
\label{sec:ledgers}

\code{stoich.py} assigns three extra rows per column from the rate's
\code{weak\_type}:

\begin{center}
\begin{tabular}{@{}lcccl@{}}
\toprule
\textbf{\code{weak\_type}} & \textbf{$d_{e^-}$} & \textbf{$d_\nu$} &
\textbf{$d_{\bar\nu}$} & \textbf{effect on \Ye} \\
\midrule
\code{electron\_capture} & $-1$ & $+1$ & 0 & lowers \\
\code{beta\_neg}         & $+1$ & 0    & $+1$ & raises \\
\code{beta\_pos}         & $-1$ & $+1$ & 0 & lowers \\
\bottomrule
\end{tabular}
\end{center}

Then $\nut = \mathrm{vstack}(\nu, d_e, d_\nu, d_{\bar\nu})$, shape
$(n{+}3) \times r$.

\paragraph{The trap, and why the docstring shouts about it.} There is an obvious
way to fill these rows: read $\Delta Z_j = \sum_i Z_i\nu_{ij}$ off the nuclear
part and set $d_e = \Delta Z$. It would always work, would never fail a test,
and would make the entire charge-to-lepton closure check \textbf{a tautology}
--- you would be verifying that a number equals itself.

Instead the ledger is assigned from the rate's \emph{declared type}, which comes
from a completely different place: pynucastro's classification of the rate
object (\code{rate.weak\_type}, ultimately from the REACLIB\,/\,tabular
metadata), reconciled against MESA in Tier~1 \S VI. The closure check

\begin{srclisting}
if dq != d_e[j]:
    raise ValueError(f"... nuclear dZ = {dq} but ledger d_electron = {d_e[j]} "
                     "- charge-to-lepton closure broken")
\end{srclisting}

is then a genuine cross-check between two independent descriptions of the same
reaction: \emph{the stoichiometry of the nuclei} and \emph{the declared weak
channel}. If a rate were mislabelled --- an EC filed as a $\beta^-$ --- this
catches it, loudly, by name, at build time.

\begin{invariant}
This is Tier~1 \S VIII.D's habit applied prospectively rather than
retrospectively: \emph{ask of any green test what would still be wrong if it
passed.} A back-derived ledger passes always and detects nothing. Two of the
build's five fail-loud validations ($|\Delta Z| \ne 1$; \code{dq != d\_e[j]})
exist only because the ledger is independent.
\end{invariant}

\paragraph{The lepton-number row.} The third row makes lepton number close per
column:
%
\begin{equation}
\Delta L_j \;=\; d_{e^-,j} + d_{\nu,j} - d_{\bar\nu,j} \;=\; 0.
\label{eq:leptonclose}
\end{equation}
%
Check each: EC gives $(-1) + (+1) - 0 = 0$~\yes; $\beta^-$ gives
$(+1) + 0 - (+1) = 0$~\yes; $\beta^+$ gives $(-1) + (+1) - 0 = 0$~\yes. The sign
on $\bar\nu$ is $-1$ because an antineutrino carries lepton number $-1$.
\code{test\_toy\_lepton\_number\_closes\_per\_weak\_column} and
\code{test\_exported\_lepton\_ledgers\_close} assert exactly this, and
additionally that the ledger rows are identically zero on strong columns.

% =====================================================================
\subsection{\texorpdfstring{$\beta^+$}{beta+} and the annihilated positron}
\label{sec:betaplus}

$\beta^+$ and EC share a ledger signature --- $(-1, +1, 0)$ --- which looks
wrong at first sight, because $\beta^+$ \emph{emits} a positron whereas EC
\emph{absorbs} an electron.

The resolution is the plasma. A positron emitted into silicon-burning matter at
$\Tn \sim 4$, $\rho \sim 10^{8}$ has a mean free path against annihilation
vastly shorter than anything else in the problem: it finds an ambient electron
and annihilates essentially instantly (the $e^+e^- \to 2\gamma$ cross-section at
thermal energies, times $n_e \sim 10^{31}$\,cm$^{-3}$, gives a lifetime far
below any timestep in the regime box). So the \emph{net} effect of a $\beta^+$
decay on the composition is:

\begin{srclisting}
  nucleus:      Z -> Z-1              (dZ_nuclear = -1)
  positron:     emitted, then annihilates with a plasma electron
  net electron ledger:  -1           (one electron removed from the sea)
  neutrino:     +1 nu_e emitted
  photons:      2 x 511 keV, thermalised
\end{srclisting}

Charge closes: the nucleus loses one unit, the electron sea loses one electron,
total change zero. Lepton number closes through the $\nu_e$: the emitted $e^+$
($L = -1$) and the absorbed $e^-$ ($L = +1$) cancel, leaving the $\nu_e$'s $+1$
balanced against\ldots{} nothing? No --- the emitted $e^+$ has $L = -1$, so the
pair $(e^+, \nu_e)$ from the decay already has $\Delta L = 0$, and the
\emph{subsequent} annihilation removes an ambient $e^-$ ($L = +1$) together with
the $e^+$ ($L = -1$), also $\Delta L = 0$. In the ledger's bookkeeping, which
tracks only the persistent species, the composite process is ``one electron
removed, one neutrino emitted'': $(-1) + (+1) - 0 = 0$.~\yes

\textbf{Both EC and $\beta^+$ therefore lower \Ye; only $\beta^-$ raises it.}
Which is why the weak-sector census counts \emph{directions}: \msn{80} has 46
weak columns $= 21$ EC $+$ 19 $\beta^-$ $+$ 6 $\beta^+$ (27 lower \Ye, 19
raise); \msn{151} has $174 = 84 + 85 + 5$ (89 lower, 85 raise).
\resultsref{2026-07-08} --- note this is the provisional 610\,/\,1522 export;
the reconciled export has 46\,/\,173 weak columns.\derivedhere

The 511\,keV photons are \emph{not} in any ledger, and that is a deliberate
omission worth noting: they go into the thermal bath, i.e. into \enuc.
\S\ref{sec:energyidentity}'s energy identity accounts for them only insofar as
the rate $Q$-value does. Tier~1 \S VIII.C.3's dropped-neutrino-column item lives
in the same neighbourhood.

% =====================================================================
\subsection{The constraint matrix \texorpdfstring{$C$}{C}}
\label{sec:cmatrix}

\begin{srclisting}
def constraint_matrix(A, Z):
    n = len(A)
    C = np.zeros((3, n + 3))
    C[0, :n] = A                # baryon number
    C[1, :n] = Z; C[1, n] = -1  # charge  (electron carries -1)
    C[2, n] = +1; C[2, n+1] = +1; C[2, n+2] = -1   # lepton number
    return C
\end{srclisting}

Three rows over the extended species list $[\text{nuclei}\ldots, e^-, \nu,
\bar\nu]$:

\begin{center}
\begin{tabularx}{\textwidth}{@{}L{1.6cm} L{3.4cm} Y@{}}
\toprule
\textbf{row} & \textbf{reads} & \textbf{why the entries are what they are} \\
\midrule
baryon & $\sum_i A_i \dd Y_i$ & leptons have $A = 0$ \\
charge & $\sum_i Z_i \dd Y_i - \dd Y_{e^-}$ &
  the electron's charge is $-1$, so its \emph{column} entry is $-1$; $\nu$,
  $\bar\nu$ are neutral \\
lepton & $\dd Y_{e^-} + \dd Y_\nu - \dd Y_{\bar\nu}$ &
  $e^-$ and $\nu$ carry $L = +1$, $\bar\nu$ carries $-1$; nuclei carry 0 \\
\bottomrule
\end{tabularx}
\end{center}

And the build asserts \code{max |C @ nu\_ext| == 0.0} exactly, for the same
integer-arithmetic reason as \S\ref{sec:exactzero}.

\paragraph{Why the charge row is the interesting one.} It does \emph{not} say
``nuclear charge is conserved''. It says nuclear charge change is exactly
matched by the electron ledger. \Ye{} is free to move; the \emph{books} must
balance. This is the formal version of \S\ref{sec:leak}'s third point.

\paragraph{$\rank(C) = 3$, and $\cond(CC\Tr) = 3.95\times10^{4}$\,/\,%
$1.03\times10^{5}$.} \derivedhere; matches \resultsref{2026-07-08}. The
conditioning is entirely a units artefact: row 0 has entries up to 151, row 2
has entries of size 1, so $CC\Tr$ has diagonal entries spanning ${\sim}4$--5
orders. \code{build\_projector} sidesteps it by never forming $CC\Tr$
(\S\ref{sec:projector}). Nothing here is near-singular in any physical sense ---
rank-deficiency would mean one of the three bookkeeping laws was a combination
of the others, which it manifestly is not, and
\code{test\_projector\_rejects\_rank\_deficient\_C} checks that the guard fires
when it is.

% =====================================================================
\subsection{Target A's theorem}
\label{sec:targetatheorem}

Now everything is in place, and the central result is one line.

\begin{invariant}
\textbf{Theorem (Target A conserves).} For \emph{any} $\netflux \in
\mathbb{R}^r$, $\dd\Yt = \nut\netflux$ satisfies
$C\cdot \dd\Yt = (C\nut)\netflux = 0\cdot\netflux = 0$.
\end{invariant}

That is it. No assumption on \netflux: not smallness, not sign, not smoothness,
not that it came from a trained model. A randomly initialised flux head
satisfies it. A flux head predicting nonsense satisfies it. A flux head with
NaNs satisfies it in the sense that the constraint residual is NaN$\cdot$0 ---
which is why the gate also checks finiteness implicitly through the magnitude
sweep.

\textbf{Accuracy determines \emph{which} conserving update is produced; it can
never determine \emph{whether} it conserves.} That sentence is the docstring of
\code{tests/test\_conservation.py} and the design principle of the entire
project.

Two immediate corollaries used later:

\paragraph{Corollary 1 (masking is safe).} If the flux head's output is
multiplied elementwise by any gate $g \in \mathbb{R}^r$, the result
$\nut(g \odot \netflux)$ is still in $\operatorname{null}(C)$ --- because
$g \odot \netflux$ is still just some vector in $\mathbb{R}^r$. Column-scaling a
stoichiometric matrix cannot break its column sums.
\code{test\_toy\_mask\_never\_touches\_weak\_columns} and
\code{test\_exported\_mask\_contract\_on\_real\_matrix} verify this with a
\emph{random} gate, which is the right test: the point is that no property of
the gate is required.

But note carefully what that corollary does \emph{not} license. Masking is safe
for conservation and unsafe for physics: gating a weak column to zero conserves
everything perfectly while deleting \dYe. Hence invariant \#2 is a
\emph{separate} structural rule --- \code{eligible\_mask} takes
\code{weak\_mask} and excludes it by construction --- rather than something
conservation could have enforced.

\paragraph{Corollary 2 (the update is exact in float64 too).}
$\dd Y = \nu\netflux$ is a matvec with integer entries; the error is bounded by
$\varepsilon\cdot(|\nu|\cdot|\netflux|)$, the gross throughput, which is exactly
the tolerance shape of \S\ref{sec:exactzero}.

% =====================================================================
\subsection{What Target A does \emph{not} guarantee --- positivity}
\label{sec:positivity}

\S\ref{sec:targetatheorem}'s theorem holds for \emph{any} \netflux, and its
corollaries drew out what that buys. This subsection draws out what it does
\textbf{not}. Conservation is enforced to machine precision by a matrix
multiply; \textbf{positivity is not enforced at all}, at any point in the
emulator design. The same ``for all \netflux'' quantifier is responsible for
both.

\subsubsection{The geometry}
\label{sec:posgeometry}

The physical state space is not a linear subspace. It is the polytope
%
\begin{equation}
\mathcal{P} \;=\; \Bigl\{\,Y \in \mathbb{R}^n \ :\ Y_i \ge 0\ \forall i,\ \
                  \textstyle\sum_i A_iY_i = 1 \,\Bigr\},
\label{eq:polytope}
\end{equation}
%
a simplex-like object with $n$ facets. In CRN language (\S\ref{sec:crn}) the
reachable set from $Y_0$ is the \textbf{stoichiometric compatibility class}
$(Y_0 + \image(\nu)) \cap \mathcal{P}$ --- Tier~0 \S IV.2's ``reachable
manifold'', now with its correct geometry: an affine subspace intersected with a
cone.

Conservation-by-construction handles the \emph{affine subspace} exactly. Nothing
handles the \emph{cone}.

\subsubsection{Why it is not a small problem here}
\label{sec:posnotsmall}

\begin{itemize}
\item \textbf{$\dd Y = \nu\netflux$ can take any species negative}, for the same
  reason it conserves: the theorem holds for arbitrary \netflux, and
  ``arbitrary'' includes fluxes that drain a species past zero. The projector
  $P$ has the same property.
\item \textbf{Abundances span 15+ decades}, so ``slightly negative'' is not
  well-defined --- a species at $10^{-25}$ is thirty orders below the dominant
  one, and an absolute error of $10^{-18}$ makes it negative.
\item \textbf{Rates are built from products of $Y$} (\S\ref{sec:lambdatoR}). A
  negative $Y$ propagates into $R$ with an unphysical sign, and for a column
  with a repeated reactant ($Y^3$ in triple-$\alpha$) it propagates with an
  unphysical \emph{magnitude} too.
\item \textbf{In rollout it compounds}: a negative abundance produces a negative
  destruction rate, which \emph{grows} the negative value. This is a divergence
  mechanism entirely separate from the usual rollout-instability story, and it
  belongs in the Tier-4 S20 node that does not exist yet.
\end{itemize}

The reference integrator does address it, minimally and correctly: $Y$ is
clipped at 0 \textbf{inside rate evaluation only}, never in the solver state
(\S\ref{sec:bdf}), so the regularisation cannot inject mass. That is the right
pattern, and it is a \emph{solver} pattern --- nothing equivalent exists on the
emulator side.

\subsubsection{What the literature offers}
\label{sec:poslit}

Two families, and the second is the one the project has not met:

\begin{enumerate}
\item \textbf{Positivity-preserving projection} --- project onto the constraint
  manifold, then backtrack along the correction until positivity holds. The
  architecture docs cite \citet{kircher2025} for exactly this --- atom-balance
  projection with linear-interpolation backtracking, in ML surrogates for
  chemical kinetics --- so this is known to the project at the \emph{citation}
  level; it is absent from the study map, from the code, and from the gate list.
\item \textbf{Modified Patankar schemes} --- \citet{burchard2003}. These are
  built for exactly the system shape here, $\dot y = P(y) - D(y)$, and are
  \textbf{unconditionally positive \emph{and} exactly conservative},
  simultaneously, for any step size. The trick is to weight \emph{both} the
  production and the destruction terms by $y^{n+1}/y^{n}$ factors, which makes
  the scheme implicit in a way that cannot cross zero --- weighting destruction
  alone (the classic \citealp{patankar1980} trick) buys positivity but breaks
  conservation; the symmetric weighting is what restores it. That is a strong
  existence result: \textbf{conservation and positivity are not in tension}, and
  a scheme achieving both already exists in the production--destruction
  literature.
\end{enumerate}

Whether an MPRK-style construction can be made differentiable and cheap enough
to sit in a decode step is an open question --- but the fact that the two
constraints are jointly achievable is exactly the sort of thing you want to know
before designing around a supposed trade-off.

\subsubsection{The minimum that should exist now}
\label{sec:posminimum}

A diagnostic, not an architecture: on any predicted update, report
\code{min\_i (Y\_i + dY\_i)} and the fraction of (state, species) cells going
negative, stratified by abundance decade. Zero cost, and it converts an unknown
into a number. Open (register: \reg{R-04}).

% =====================================================================
\subsection{The Target B projector, derived twice}
\label{sec:projector}

Target B predicts $\dd\Yt$ directly and repairs it. The repair operator:

\subsubsection{Derivation 1 --- minimum correction (Lagrange)}
\label{sec:projlagrange}

Given a raw output $v \in \mathbb{R}^{n+3}$, find the nearest point on the
constraint manifold:
%
\begin{equation}
\min_{w}\ \tfrac12\|w - v\|_2^2 \quad\text{s.t.}\quad Cw = 0 .
\label{eq:projmin}
\end{equation}
%
Lagrangian $L = \tfrac12\|w-v\|^2 + \lambda\Tr Cw$. Stationarity:
$w - v + C\Tr\lambda = 0 \Rightarrow w = v - C\Tr\lambda$. Impose the
constraint: $Cv - CC\Tr\lambda = 0 \Rightarrow \lambda = (CC\Tr)^{-1}Cv$
(invertible because $\rank(C) = 3$). Substitute:
%
\begin{equation}
w \;=\; \bigl(I - C\Tr(CC\Tr)^{-1}C\bigr)\,v \;\equiv\; P v .
\label{eq:projector}
\end{equation}
%
So \textbf{$P$ is the orthogonal projector, and orthogonality is a choice}: it
is the projector that makes the \emph{smallest Euclidean correction}. That is a
real modelling assumption, not a mathematical necessity --- any oblique
projector onto $\operatorname{null}(C)$ would also conserve. The Euclidean
metric treats a correction of $10^{-6}$ to $Y_{\rm neut}$ and to $Y_{\rm si28}$
as equally costly, which is not obviously right when their abundances differ by
ten orders of magnitude. See \reg{R-17}.

\subsubsection{Derivation 2 --- QR, which is what ships}
\label{sec:projqr}

\begin{srclisting}
Q, R = np.linalg.qr(C.T)      # reduced: Q is (m, 3) with orthonormal columns
P = np.eye(m) - Q @ Q.T
\end{srclisting}

$Q$'s columns are an orthonormal basis for
$\mathrm{rowspace}(C) = \mathrm{colspace}(C\Tr)$. $QQ\Tr$ is the orthogonal
projector \emph{onto} that space, so $I - QQ\Tr$ projects onto its orthogonal
complement, which is $\operatorname{null}(C)$. Equality with Derivation 1: write
$C = R\Tr Q\Tr$. Then $CC\Tr = R\Tr Q\Tr QR = R\Tr R$, and
%
\begin{equation}
C\Tr(CC\Tr)^{-1}C \;=\; QR(R\Tr R)^{-1}R\Tr Q\Tr
\;=\; QRR^{-1}R^{-\mathsf{T}}R\Tr Q\Tr \;=\; QQ\Tr. \qquad\yes
\label{eq:qrequiv}
\end{equation}

\paragraph{Why QR is the one that ships.} The normal-equations form requires
inverting $CC\Tr$, whose condition number is $\cond(C)^2$. With
$\cond(CC\Tr) \approx 4\times10^{4}$, $\cond(C) \approx 200$ --- so the
normal-equations route throws away ${\sim}2.3$ decimal digits for nothing. Here
that is harmless ($10^{-16} \to 10^{-14}$, still far inside the $10^{-12}$
gate), but the QR form is free and additionally gives symmetry and idempotence
to near machine precision by construction rather than by luck. The docstring
says exactly this, and the measured result confirms it:
$\|P^2 - P\|_\infty \le 1.0\times10^{-15}$, $P$ symmetric to 0.0 exactly,
$\rank(P) = m - 3$ exactly, max relative constraint residual over 20 decades of
$\dd Y = 6.68\times10^{-17}$ (\msn{80})\,/\,$1.04\times10^{-16}$ (\msn{151})
against a $10^{-12}$ gate. \resultsref{2026-07-08}

\code{build\_projector} also guards rank explicitly:

\begin{srclisting}
if np.abs(np.diag(R)).min() <= 1e-12 * np.abs(np.diag(R)).max():
    raise ValueError("C is rank-deficient - constraint rows are not independent")
\end{srclisting}

which is a rank test on $R$'s diagonal --- valid because for a full-rank matrix
QR gives $|R_{ii}| > 0$, and the ratio to the largest is the natural scale-free
threshold.

% =====================================================================
\subsection{Why the projector acts on the extended vector --- the laundering
            trap}
\label{sec:laundering}

$P$ is $(n{+}3) \times (n{+}3)$, not $n \times n$. If instead you built a
nuclei-only constraint matrix $C_{\rm nuc} = [A; Z]$ ($2 \times n$) and
projected $\dd Y_{\rm nuclei}$, you would be imposing
%
\begin{equation}
\sum_i Z_i\,\dd Y_i = 0 \quad\Longleftrightarrow\quad \dd Y_e = 0 .
\label{eq:launder}
\end{equation}

\begin{warn}
\warnglyph{}~\textbf{The projector would delete the weak signal.} Not degrade it
--- delete it, exactly, by orthogonal projection. And it would do so while
making every conservation diagnostic look perfect. This is the sharpest instance
in the repo of ``ask what would still be wrong if the test passed'': a
nuclei-only projector passes baryon conservation, passes charge conservation,
passes idempotence, passes symmetry, and has destroyed the physics.
\end{warn}

\code{test\_projection\_preserves\_weak\_dYe\_signal} is the guard: build a
weak-only flux, compute $\dd\Yt = \nut\netflux$, record
$\dYe = Z\cdot \dd Y_{\rm nuclei} \ne 0$, project, and require the value to be
unchanged to $10^{-12}$ relative. Measured change: $\le 3.4\times10^{-21}$.
\resultsref{2026-07-08}

Why is it unchanged? Because $\nut\netflux$ is \emph{already} in
$\operatorname{null}(C)$ (\S\ref{sec:targetatheorem}), and $P$ is the identity on
$\operatorname{null}(C)$ --- \code{test\_on\_manifold\_input\_is\_unchanged}
states this directly. The extended projector's action on a genuinely conserving
update is the identity; only the \emph{nuclei-only} projector would move it, and
it would move it in exactly the \dYe{} direction.

\paragraph{The other half of invariant \#4: $P$ is valid in a LINEAR output
space only.} Suppose the head emits $u = \asinh(\dd Y/s)$ or a signed-log, and
you project $u$. Then $C\cdot P(u) = 0$ says a linear constraint holds \emph{in
the warped coordinates}. Mapping back,
$\sum_i A_i \cdot \sinh(u_i)\cdot s \ne 0$ in general --- the constraint you
imposed is not the constraint you wanted. This mathematical argument is the
project's own \derivedhere; what NuGNN itself documents is the empirical half:
\citet{kim2026} emit $\mathrm{sign}(f)\cdot\log_{10}|f|$ outputs, attempted
explicit conservation enforcement, found that back-transforming through the
signed log made optimisation unstable, and shipped without enforcement ---
relying on implicitly learned conservation plus a rollout-time retry heuristic
(a step whose predicted $\sum\Delta X$ drifts far from zero is retried at
smaller $\Delta t$). The rule that falls out is stated in \code{CLAUDE.md}
invariant \#4 and in \code{projector.py}'s docstring: \textbf{conservation lives
in the decode step; no nonlinear transform may sit between the conservation map
and the output.} $\asinh$\,/\,signed-log are legitimate \emph{latent}
representations (\foreignreg{0}{R-18}) --- they must simply not be the space
where the sum is evaluated.

% =====================================================================
\subsection{Target A and Target B reach the same set}
\label{sec:equivalence}

A result I had assumed was false until I measured the ranks, and which changes
how the Target A/B decision should be framed. \S\ref{sec:reachable} asserted it;
here it is proved.

\begin{invariant}
\textbf{Theorem}\derivedhere. $\image(\nut) = \operatorname{null}(C)$.
\end{invariant}

\emph{Proof.} ($\subseteq$) $C\nut = 0$, so every $\nut\netflux$ lies in
$\operatorname{null}(C)$. ($\supseteq$) Dimension count:
$\dim\operatorname{null}(C) = (n{+}3) - \rank(C) = (n{+}3) - 3 = n$. And
$\rank(\nut) = n$ --- measured: 80 for \msn{80}, 151 for \msn{151}
(\S\ref{sec:rank}). A subspace contained in another of the same finite dimension
is equal to it. $\square$

The measurement is what makes the second half non-trivial: $\rank(\nut)$
\emph{could} have been less than $n$ (a network too small to reach every
conserving direction), and for a sufficiently truncated network it would be. It
is not, for either network here.

\paragraph{What this means.}

\begin{center}
\begin{tabularx}{\textwidth}{@{}L{4.2cm} Y Y@{}}
\toprule
 & \textbf{Target A} & \textbf{Target B} \\
\midrule
output & $\netflux \in \mathbb{R}^r$ (607\,/\,1518) &
  $\dd\Yt \in \mathbb{R}^{n+3}$ (83\,/\,154) \\
conserving? & by construction ($C\nut = 0$) & by construction ($P$) \\
\textbf{reachable set} & $\mathbf{\operatorname{null}(C)}$ &
  $\mathbf{\operatorname{null}(C)}$ \\
parameterisation & redundant: 528-\,/\,1368-dim fibre & minimal \\
error amplification into $\dd Y$ & $\cond(\nut) \approx 41$\,/\,55 &
  1 ($P$ is an orthogonal projection) \\
per-output physical meaning & a named reaction's flux &
  a species' abundance change \\
maskable? & yes, per reaction & no natural analogue \\
needs \Phifl{} labels? & \textbf{yes} (\S\ref{sec:augmented}) &
  no --- $\Delta X$ suffices \\
\bottomrule
\end{tabularx}
\end{center}

So the choice is \textbf{not} about expressivity. Both can express exactly the
conserving updates and nothing else. The trade is:

\begin{itemize}
\item Target A buys physical structure --- per-reaction attribution, the
  equilibrium mask, reaction-level loss weighting, direct comparability to the
  bridge sets --- at the cost of a badly over-parameterised output (528
  unconstrained directions), a conditioning factor of ${\sim}40$--60 from
  \netflux{} to $\dd Y$, and a label requirement the shipped dataset does not
  meet.
\item Target B buys a minimal, well-conditioned parameterisation and works with
  the labels that exist, at the cost of throwing away every per-reaction handle.
\end{itemize}

Framed that way, \code{CLAUDE.md}'s switch condition --- ``Target A $\to$ Target
B if $\cond(\Sact) > 10^{6}$, or if $>{\sim}90\%$ of reactions carry net flux
below the \Ye{} floor'' --- reads correctly: it is triggered by the
\emph{conditioning} and the \emph{structure}, which is exactly what Target A is
being bought for. When the structure is not there, Target A's costs remain and
its benefits do not.

That reframing has a documentation consequence, which is why this section is a
register anchor: no project document should frame the Target A vs B choice as an
\emph{expressivity} difference, because the theorem above says it is not. Open
(register: \reg{R-22}).

% =====================================================================
\subsection{The effective dimension of the reachable set --- and where
            \texorpdfstring{\Ye}{Ye} hides}
\label{sec:effdim}

\S\ref{sec:equivalence} established \emph{which increments} are reachable --- the
whole of $\operatorname{null}(C)$, by either target. The complementary question
is which \textbf{states} are actually visited. Tier~0 \S IV.2 introduced the
reachable manifold qualitatively; nothing anywhere puts a number on it, and the
number matters in a way I did not anticipate.

\subsubsection{The measurement: how many directions the compositions occupy}
\label{sec:effdim-measurement}

PCA over pre-stall rows of 80 randomly chosen shipped trajectories per network
($\le 25$ log-spaced rows each;
\code{data/trajectories.select\_rows(prestall=True)}). \derivedhere:

\begin{srclisting}
X = np.vstack(...)                     # (rows, n_species) mass fractions
Xc = X - X.mean(0)
s = np.linalg.svd(Xc, compute_uv=False)
c = np.cumsum(s**2) / (s**2).sum()
\end{srclisting}

\begin{center}
\begin{tabularx}{\textwidth}{@{}Y L{3.2cm} L{3.2cm}@{}}
\toprule
 & \textbf{\msn{80}} ($n = 80$) & \textbf{\msn{151}} ($n = 151$) \\
\midrule
rows\,/\,trajectories & 2079\,/\,80 & 2076\,/\,80 \\
PCs for 90\% of $\Var(X)$ & \textbf{7} & \textbf{14} \\
PCs for 99\% & 22 & 38 \\
PCs for 99.9\% & 37 & 61 \\
participation ratio $(\sum\lambda)^2/\sum\lambda^2$ & 3.42 & 8.58 \\
PCs for 90\% of $\Var(X)$ in \textbf{$\log_{10}$} space & 6 & 5 \\
\bottomrule
\end{tabularx}
\end{center}

\textbf{The compositions actually visited occupy an effective manifold of
dimension $\approx 3$--14, not 80 or 151.} The participation ratio --- a
basis-free effective rank --- is 3.4 and 8.6.

This is worth internalising, because it reframes the difficulty of the learning
problem. It is very likely \emph{why} a plain feed-forward network (the NNN
baseline, S16) works at all on a nominally 151-dimensional map: the target is
not a 151-dimensional function, it is a ${\sim}10$-dimensional one embedded in
151 coordinates. It also gives the size-transfer question (\msn{80} $\to$
\msn{151}) a sharper form: the effective dimension roughly \emph{doubles}
($3.4 \to 8.6$), which is a more informative statement about ``how much harder
is the big network'' than $80 \to 151$ is.

\begin{warn}
\textbf{Caveats, stated up front.} These are the \emph{shipped} trajectories,
which stall at a bug-displaced attractor (Tier~3, S13) and are
constant-$(T, \rho)$; 80 files per network; linear-$X$ PCA is dominated by the
few abundant species. The number is a lower bound on the diversity a full corpus
would show, and the measurement should be repeated on the clean reruns and
across the $(T, \rho)$ box. The \emph{qualitative} conclusion --- effective
dimension $\ll n$ --- is robust; the specific integer is not.
\end{warn}

\subsubsection{The part that changes a design decision}
\label{sec:effdim-design}

The obvious next move is dangerous, so it is worth doing the second measurement
before anyone makes it: if the manifold is ${\sim}10$-dimensional, why not
compress --- a PCA basis, a latent bottleneck, a variance-weighted loss?

Because \textbf{\Ye{} does not live in the leading directions.}
$\Ye = \sum_i (Z_i/A_i)X_i$ is a linear functional $w\cdot X$, so the fraction
of $\Var(\Ye)$ captured by the top-$k$ PCs of $X$ is computable exactly.
\derivedhere:

\begin{center}
\begin{tabularx}{\textwidth}{@{}Y L{3.2cm} L{3.2cm}@{}}
\toprule
 & \textbf{\msn{80}} & \textbf{\msn{151}} \\
\midrule
std$(\Ye)$ across the sample & $1.18\times10^{-2}$ & $1.61\times10^{-2}$ \\
PCs for 90\% of \textbf{$\Var(X)$} & 7 & 14 \\
PCs for 90\% of \textbf{$\Var(\Ye)$} & \textbf{16} & 11 \\
PCs for 99\% of $\Var(X)$ & 22 & 38 \\
PCs for 99\% of \textbf{$\Var(\Ye)$} & \textbf{43} & \textbf{65} \\
cumulative $\Var(X)$\,/\,$\Var(\Ye)$ at $k = 10$ &
  0.9491\,/\,0.8716 & 0.8565\,/\,0.8799 \\
\ldots at $k = 20$ & 0.9875\,/\,0.9082 & 0.9521\,/\,0.9426 \\
\ldots at $k = 30$ & \textbf{0.9966\,/\,0.9295} &
  \textbf{0.9811\,/\,0.9527} \\
\bottomrule
\end{tabularx}
\end{center}

\begin{warn}
\warnglyph{}~Read this carefully, because the effect is \textbf{in the tail, not
in the leading directions}. At small $k$ the two are comparable --- \msn{151} at
$k = 10$ even captures \Ye{} slightly \emph{better} than $X$. What differs is
convergence: to reach 99\%, \Ye{} needs \textbf{43 vs 22} components (\msn{80})
and \textbf{65 vs 38} (\msn{151}), i.e. 1.7--2$\times$ as many; and by $k = 30$,
$X$ has converged to 99.7\%\,/\,98.1\% while \Ye{} has stalled at
93.0\%\,/\,95.3\%. \textbf{The last few percent of the \Ye{} signal is smeared
across many low-variance directions and cannot be recovered by truncation.}
\end{warn}

\subsubsection{The arithmetic that settles it}
\label{sec:effdim-arithmetic}

That tail behaviour would be a mild caution if the \Ye{} tolerance were loose.
It is not. The per-step budget is $3\times10^{-6}$ (\S\ref{sec:budget}) against
$\mathrm{std}(\Ye) = 1.18\times10^{-2}$, so the tolerable \emph{relative} error
is
%
\begin{equation}
\frac{3\times10^{-6}}{1.18\times10^{-2}} \;=\; 2.5\times10^{-4},
\label{eq:reltol}
\end{equation}
%
and a rank-$k$ truncation's \Ye{} error is
$\sqrt{1 - \mathrm{cumVar}_{\Ye}(k)}\cdot\mathrm{std}(\Ye)$. Meeting the budget
therefore requires
%
\begin{equation}
\mathrm{cumVar}_{\Ye}(k) \;\ge\; 1 - (2.5\times10^{-4})^2
\;=\; 1 - 6.5\times10^{-8} \;=\; 99.999993\% .
\label{eq:cumvarreq}
\end{equation}
\derivedhere

\textbf{No truncation achieves that.} Not $k = 30$, not $k = 60$, not any $k$
short of the full rank --- the measured curve is at 93\% by $k = 30$. So this is
not ``variance compression is risky for \Ye''; it is \textbf{arithmetically
excluded}: any variance-truncated reduction of the composition --- PCA basis,
POD, a linear autoencoder bottleneck sized by reconstruction error --- destroys
the \Ye{} budget by three to four orders of magnitude before the model has made
a single error of its own.

\subsubsection{And it is not an accident --- it is the rank measurement}
\label{sec:effdim-stoich}

This is not an empirical quirk of these trajectories; it is forced by the
stoichiometry, and the proof is a result already in this document.

\S\ref{sec:rank} measured that the left null space of $\nu$ restricted to
strong\,/\,EM columns is exactly $\mathrm{span}\{A, Z\}$. So \emph{every}
increment produced by the strong sector satisfies $Z\cdot \dd Y = 0$, i.e.
$\dYe = 0$ exactly. The strong sector is fast, carries essentially all the
composition change, and is therefore responsible for essentially all the
composition variance --- and every bit of it is \Ye-neutral.

\begin{result}
\textbf{\Ye{} moves only through the weak sector, which is slow. The dominant
variance directions are therefore \Ye-blind by construction, and the \Ye{}
signal is pushed into the low-variance tail.}\derivedhere
\end{result}

Consequences, all actionable:

\begin{enumerate}
\item \textbf{Never size a latent bottleneck by explained composition variance.}
  The requirement is $1 - 6.5\times10^{-8}$ of $\Var(\Ye)$; reconstruction error
  is simply the wrong instrument, by four orders of magnitude. If a bottleneck
  is used at all, \Ye{} must be carried \emph{outside} it --- as an explicit
  conserved coordinate, not as something reconstructed.
\item \textbf{A plain MSE-on-$X$ (or on $\Delta X$) loss is variance-weighted},
  hence subject to the same blindness. An explicit \Ye{} term is not a nicety,
  it is structurally required. (\code{CLAUDE.md}'s loss-weighting discussion and
  the $|\dot{\dYe}|$ concentration measurements --- top-5 weak channels carrying
  0.96\,/\,0.75, \resultsref{2026-07-12} --- are the right instinct; this is the
  quantitative argument for it.)
\item \textbf{Any PCA\,/\,POD\,/\,autoencoder-based reduction must be validated
  on \Ye{} separately}, not on reconstruction error.
\item \textbf{It suggests a genuinely better decomposition}: split the update as
  $\dd Y = \dd Y_{\rm strong} + \dd Y_{\rm weak}$ with
  $\dd Y_{\rm strong} \in \operatorname{null}(Z\Tr)$ enforced structurally. The
  strong part then lives in a lower-dimensional, well-conditioned space, and the
  weak part --- small, slow, and \emph{the entire \Ye{} signal} --- is a
  separate head with its own loss and its own error budget. That is a Tier-4
  architecture idea and, like \S\ref{sec:entropyactive}'s, it needs a novelty
  check before anyone gets attached to it.
\end{enumerate}

Open (register: \reg{R-02}) --- repeat on the clean reruns and across the
$(T, \rho)$ box before the Tier-4 choices above are made.

% =====================================================================
\subsection{The bipartite graph, the radius, and \texorpdfstring{$K$}{K}}
\label{sec:graphK}

\code{export.bipartite\_digraph} builds a heterogeneous digraph: isotope nodes
and reaction nodes; I$\to$R edges for reactants (coefficient = consumed count,
positive), R$\to$I edges for products (signed net production). A catalyst-like
species gets both.

\code{metrics.graph\_extents} reports three undirected views.
\resultsref{2026-07-08}, identical for both networks:

\begin{center}
\begin{tabularx}{\textwidth}{@{}L{5.2cm} L{1.6cm} L{1.9cm} Y@{}}
\toprule
\textbf{view} & \textbf{radius} & \textbf{diameter} & \textbf{reading} \\
\midrule
bipartite (incidence only) & 3 & 6 &
  I$\to$R$\to$I = 2 hops, so this counts half-reactions \\
bipartite $+$ I$\to$I coupling edges & 2 & 4 & shortcut edges added \\
isotope projection (I$\to$I only) & \textbf{1} & \textbf{2} &
  one reaction = one hop \\
\bottomrule
\end{tabularx}
\end{center}

All connected, one component.

\textbf{Isotope-projection radius 1} is the hub structure of
\S\ref{sec:numap}(c) restated: there exists a species (\nuc{He}{4}) adjacent to
\emph{every} other species --- it shares a reaction with all of them. Diameter
2: any two species are at most two reactions apart.

\textbf{$K = \lceil\text{radius}\rceil + 2 = 5$} on the bipartite graph is the
checklist's implied processor depth (\code{GraphExtent.implied\_K}). The
reasoning that makes it coherent: in the bipartite view one \emph{reaction hop}
costs two message-passing layers, so $K = 5$ layers $\approx 2.5$ reaction hops
--- which covers the isotope-projection \emph{diameter} of 2 reaction hops with
one layer to spare. That is the honest justification; note that the formula as
written uses \textbf{radius}, which for a graph-level readout from a central
node would be the right base, but for a node-level output (which this is ---
every species needs its own $\dd Y$) the stricter base is eccentricity, and the
fact that $K = 5$ also clears the diameter is a happy coincidence of this
particular graph rather than something the formula guarantees. Logged as an open
item in \reg{R-19}.

\textbf{The hazard the number hides.} A radius of 1 in the isotope projection
means the graph carries almost no long-range structure to learn: everything is
two hops from everything through three hub nodes. Depth beyond $K \approx 4$--5
buys nothing topologically, and --- more importantly --- an architecture that
appears to ``use the graph'' may in fact be routing all information through
$n$\,/\,$p$\,/\,\nuc{He}{4} node features. That is a Tier-4
experimental-design problem, but it is \emph{this} number that raises it.

% =====================================================================
\subsection{Code walkthrough}
\label{sec:codes7}

Reading order, ${\sim}600$ lines total.

\paragraph{\code{graph/isotopes.py}} --- the isotope table: names, $A$, $Z$,
pynucastro \code{Nucleus} objects, and \code{index()}. The row order of $\nu$ is
the YAML order, which is the CSV-header order of the training data. This is the
join key for everything. (Tier~0 \S II.7 covers it.)

\paragraph{\code{graph/network.py}} --- builds the pynucastro
\code{RateCollection} with duplicate resolution and the reconciliation
disposition applied. Tier~1 \S VI.4.

\paragraph{\code{graph/stoich.py}} (211 lines) --- the node. Read in this order:
\begin{enumerate}
\item The module docstring --- it \emph{is} the spec for the conservation layer.
\item \code{WEAK\_LEDGERS} --- the three-row table of \S\ref{sec:ledgers}.
\item \code{constraint\_matrix} --- \S\ref{sec:cmatrix}.
\item \code{build\_stoich}: the $\nu$ loop, then the \textbf{five fail-loud
  validations} (baryon leak; strong-column charge leak; weak $|\Delta Z| \ne 1$;
  ledger\,/\,$\Delta Z$ mismatch; unrecognised weak classification), then
  \code{C @ nu\_ext} residual, then the frozen dataclass.
\item \code{Stoich} --- note what it carries beyond $\nu$: \code{weak\_mask},
  \code{weak\_type}, \code{Q}, \code{is\_tabular}, \code{derived\_from\_inverse},
  \code{chapter}, \code{source\_label}. Every one of those is a Tier-1 output
  becoming a Tier-2 column attribute (Tier~1's ``threads that carry forward'').
\end{enumerate}

\paragraph{\code{graph/projector.py}} (47 lines) --- \S\ref{sec:projector}. Two
functions; the docstring carries the invariant.

\paragraph{\code{graph/export.py}} --- serialisation. Two details worth
noticing: \code{content\_hash} hashes the \emph{content} (sorted keys, dtypes,
shapes, bytes) rather than the file, because \code{.npz} embeds zip timestamps
and would otherwise differ between identical exports. This is Tier~0 \S VII.5
(non-regenerable artifacts are persisted) meeting reproducibility: the artifact
\emph{is} regenerable, so what gets recorded in \code{RESULTS.md} is a content
hash that proves the regeneration is identical.

\paragraph{\code{graph/metrics.py}} --- the measurement layer:
\code{reaction\_census}, \code{weak\_inventory}, \code{graph\_extents},
\code{condition\_numbers}, \code{drift\_metrics}. Everything here exists to
produce a \code{RESULTS.md} row. Note \code{\_SVD\_RANK\_RTOL = 1e-12} and the
\emph{rank-revealing} condition number (ratio of extreme \textbf{nonzero}
singular values): using the naive $\sigma_{\max}/\sigma_{\min}$ would report
$\infty$, because the structural left-null vector $A$ guarantees a zero singular
value. The same definition is reused by \code{killtest.active\_set.cond\_s\_active},
and that consistency is deliberate --- the kill-test's cond gate would be
meaningless if the two definitions differed.

\paragraph{\code{scripts/export\_stoich\_matrix.py}},
\textbf{\code{scripts/check\_conservation.py}},
\textbf{\code{scripts/check\_projector.py}} --- the reproducible entry points
behind the \code{RESULTS.md} rows.

% =====================================================================
\subsection{The oracle: what the 30-test gate actually asserts}
\label{sec:oracles7}

\code{tests/test\_conservation.py} has two layers, and the \emph{layering} is
the design.

\paragraph{Layer 1 --- the synthetic toy network.} Ten species, eight columns,
hand-built in the test file: triple-$\alpha$, an $\alpha$-capture chain, an
$(\alpha,\gamma)$, a $(p,\gamma)$ bottleneck, a photodisintegration, and
\textbf{one of each weak type}. No pynucastro. No I/O.

Its purpose is to isolate \emph{the convention} from \emph{the export}. If Layer
1 fails, the bookkeeping convention itself is wrong --- the ledger table, the
sign of $\bar\nu$ in the lepton row, the $\beta^+$ treatment. If Layer 1 passes
and Layer 2 fails, the convention is fine and the pynucastro export is wrong.
Two failure modes, cleanly separated by construction. This is worth stealing as
a habit.

Layer 1 also encodes the physics-direction assertions that the production
matrices cannot conveniently state:

\begin{srclisting}
@pytest.mark.parametrize(("col", "expected_sign"),
    [(5, -1.0), (6, +1.0), (7, -1.0)],
    ids=["EC_lowers_Ye", "beta_minus_raises_Ye", "beta_plus_lowers_Ye"])
\end{srclisting}

--- i.e. \S\ref{sec:betaplus}'s conclusion, as a test with human-readable ids.

\paragraph{Layer 2 --- the production matrices.} Auto-generated by the session
fixture in \code{conftest.py} if absent (${\sim}7$\,s/net), so this layer
\textbf{never skips}. That is a deliberate anti-pattern-avoidance: a gate that
silently skips when its data is missing is not a gate. Assertions:

\begin{center}
\begin{tabularx}{\textwidth}{@{}L{4.5cm} Y L{2.2cm}@{}}
\toprule
\textbf{test} & \textbf{asserts} & \textbf{tolerance} \\
\midrule
\code{test\_exported\_shapes\_and\_species} &
  $n = 80/151$, $C$ is $3 \times (n{+}3)$, $\nut$ is $(n{+}3) \times r$ &
  exact \\
\code{test\_exported\_columns\_conserve\_exactly} &
  per column: $\sum A\nu = 0$; strong $\sum Z\nu = 0$; weak $|\Delta Z| = 1$ and
  $\Delta Z = d_e$ & \textbf{\code{== 0.0}} (\S\ref{sec:exactzero}) \\
\code{test\_exported\_lepton\_ledgers\_close} &
  $\Delta L = 0$ on weak; ledgers zero on strong; $\max|C\nut| = 0$ &
  \code{== 0.0} \\
\code{test\_exported\_weak\_mask\_matches\_weak\_type} &
  \code{weak\_mask} $\Longleftrightarrow$ \code{weak\_type != ""}; types
  $\subseteq$ the three known & exact \\
\code{test\_exported\_random\_flux\_drift} &
  \textbf{the blocking gate}: 10 seeds $\times$ 6 scales
  (\code{1e0}\ldots\code{1e-20}) $\times$ 2 networks; baryon drift and
  charge-to-lepton closure within
  $\max(10^{-12}s,\ 10^{-13}G)$; \dYe{} through weak columns $\ne 0$ & scaled \\
\code{test\_exported\_mask\_contract\_on\_real\_matrix} &
  random gate on the eligible set ($=$ \textasciitilde weak,
  \emph{structurally}) leaves conservation intact and weak columns unscaled &
  scaled \\
\bottomrule
\end{tabularx}
\end{center}

The magnitude sweep is the part I would have under-specified if writing it
myself. One scale is not a test: at $\netflux \sim 1$ an absolute $10^{-12}$
bound is meaningful; at $\netflux \sim 10^{-20}$ it is vacuous, and a broken
implementation would pass. Six scales spanning 20 decades with a tolerance that
\emph{changes character} across them (absolute at the top, relative-to-gross at
the bottom) is what makes the gate informative everywhere.

\code{tests/test\_projector.py} mirrors the structure for Target B: shape,
symmetry ($\le 10^{-14}$), idempotence ($\le 10^{-13}$), rank $= m - 3$ via
eigenvalue sum, constraint residuals across the same 20 decades with a
\emph{relative} bound ($P$ is linear hence exactly scale-equivariant, so the
meaningful statement is relative), fixed-point on the manifold, weak-\dYe{}
preservation, and the rank-deficiency guard.

\begin{repopointer}
\textbf{\textcircled{4} SEE:} notebook \code{03-graph-conservation.ipynb} ---
Figure 1 the $\nu$ sparsity structure, Figure 2 drift vs \netflux{} magnitude
(the tolerance's two regimes visible as a kink), Figure 3 the projector residual
across 20 decades, Figure 4 the graph extents and $K$.
\end{repopointer}

% =====================================================================
\subsection{Self-check}
\label{sec:selfcheck-s7}

\begin{selfcheck}
\item State the theorem relating conserved quantities to $\nu$, with the
  quantifier.
\item left-nullity$(\nu) = 1$, left-nullity$(\nu$ restricted to strong
  columns$) = 2$, left-nullity$(\nut) = 3$. Say what each of the three numbers
  rules out.
\item Why can \code{test\_exported\_columns\_conserve\_exactly} assert
  \code{== 0.0} while \code{test\_exported\_random\_flux\_drift} cannot?
\item Why is the lepton ledger assigned from \code{weak\_type} rather than from
  $Z\cdot\nu$? What would the closure test detect if it were back-derived?
\item Derive $P = I - C\Tr(CC\Tr)^{-1}C$ from a constrained least-squares
  problem, then show the QR form is equal to it.
\item What exactly goes wrong if $P$ is built from a nuclei-only $C$? Which test
  catches it?
\item Prove $\image(\nut) = \operatorname{null}(C)$, and say which half of the
  proof required a measurement.
\item The right-nullity of $\nu$ is 528. Give two distinct consequences for the
  emulator's design.
\item Why is \code{cond\_s\_active} defined with the ratio of extreme
  \emph{nonzero} singular values rather than
  $\sigma_{\max}/\sigma_{\min}$?
\end{selfcheck}
